\scanpage{4}
On notera que les conditions $1^\circ$ à $8^\circ$ sont stables par tte extension de la base. Sauf $2^\circ$, $5^\circ$, $6^\circ$, ce sont des conditions \uline{absolues} sur $F$.

Pour prouver la représentabilité de $F$, on commencera par le faire pour $n=\dim S<+\infty$, \deleted{\unsure{en suppos}}

\textcircled{1} $\dim S=0$. On est ramené au cas $S$ local artinien, on construit un préschém. loc. qu. fini sur $S$ grâce à \textcircled{$4^\circ$}, \textcircled{$5^\circ$}, on voit qu'il représente $F$ grâce à \textcircled{$1^\circ$}a, \textcircled{$2^\circ$}, \textcircled{$3^\circ$} [cf. notes jointes, cor $2$ à prop. $1$] OK. \uline{NB} On n'a eu besoin que de $1^\circ$ à $5^\circ$

\textcircled{2} $\dim S=n>0$, $S$ \uline{local complet}. Soit $s$ son pt fermé, \deleted{\unsure{---}} et $U=S-s$, alors par l'hypothèse de récurrence [où \textcircled{$1^\circ$} : \textcircled{$5^\circ$} suffisent si $S$ de dim $1$] on sait que $F_U$ est représentable \inserted{par un $X_U$}. De plus, grâce à \textcircled{$\unsure{4^\circ|2^\circ}$}, \textcircled{$5^\circ$}, on construit une famille de schémas \uline{finis} $X_i$ sur $S$, \inserted{(utilisant $S$ complet)} \unsure{correspondants} aux pts de $F_0$, et \deleted{grâce à} grâce à \textcircled{$3^\circ$} des hom $u_i:X_i\longrightarrow F$. En vertu de la prop. $2$ des notes, on est ramené à vérifier que $u_{iU}:X_{iU}$\inserted{$\longrightarrow X_U$} \uline{immersions} \uline{ouvertes}, à images deux à deux disjointes. Ce dernier fait provient de l'\uline{hypothèse}
\margin{\uline{N.B.} Si $\dim S=1$, $S$ local (\unsure{compl.}), l'hypothèse $8^\circ$ est inutile.}
\mspage{3}
